\documentclass[10,letterpaper,roman]{article}
\begin{document}
\title{Summer Funding Plan}
\author{Alexander Marshall}
\date{February 27, 2026}
%\maketitle
I intend to use the Hamilton Award Funding to fund my summer reserch, so that I can advance towards the completion of my doctorate. In particular, I will use the funding provided by the Hamilton Award primarily to avoid having to work as a TA over the summer. In my case, this is particularly relevant, as I am trying to complete my degree before my Roster Doctoral Fellowship is completed. Doing so will allow me to foucs the entirity of my energies to reserch, speeding progress and enabling me to progress towards my research goals more quickly.\\
I am investigating the impact of strains on diffusive transport in poroelastic media.\\

Abstract goes here.\\

Over the course of the summer, I will conduct strains on cells to understand how strains impact the transport of transcription factors \textit{in situ}. This will entail adapting an extant protocol used in the research described above to act upon living cells in a controlled enviroment, which I anticipate will proceed as follows. I will cultivate a cell line, most likely human macrophages due to extensive familiarity with these cells due to prior lab experiance. I will treat them with a non-toxic HALO-tag solution to fluorescently label thier transcription factors (TFs).\\

This step will likely be done in collaboration with the lab of , who freqently uses this technique to prepare specimens for projects we collaborate on; their experiance in this matter will be invalueble.

Using the highly inclined laminated optical sheet microscopy system I built in the course of previous work, I have demonstrated the ability to conduct high-speed tracking of single nanoscopic flourophores, so this task should be well within the demonstrated technical capabilites of my current experimental setup. Additionally, a similar experiment has been conducted using HALO-tagged histone protein complexes in collaboration with Dr. Kerry Bloom of the UNC Biology Department. I anticipate the results from this experiment will be similar with regards to experimental methods, data collection, and analytical processing.  Dr. Bloom's lab freqently uses this technique to prepare specimens for projects we collaborate on; their experiance in this matter will be invalueble.\\

Our microscope operates with low phototoxicity inside a climate-controlled incubation system, allowing cell health to be maintained for extended experimental sessions. If warrented, a unique in-house instrument which enables simultaneous atomic force microscopy and light sheet fluorecence microscopy could also be used to produce simultaneous particle tracking and rheological data to examine diffusive dynamics as a function of strain.

\end{document}